\documentclass[11pt, oneside, noexaminfo]{sapthesis}   	% use "amsart" instead of "article" for AMSLaTeX format
\usepackage[margin=1.5cm]{geometry}                		% See geometry.pdf to learn the layout options. There are lots.
\geometry{letterpaper}                   		% ... or a4paper or a5paper or ... 
%\geometry{landscape}                		% Activate for rotated page geometry
%\usepackage[parfill]{parskip}    		% Activate to begin paragraphs with an empty line\usepackage{graphicx}				% Use pdf, png, jpg, or eps§ with pdflatex; use eps in DVI mode
								% TeX will automatically convert eps --> pdf in pdflatex
								
%\usepackage{{Gratzer-Color-Scheme}	% Active to color theorems red and lemmas blue
	
\usepackage{amssymb}
\usepackage[italian]{babel}			% Imposta la lingua italiana (necessaria a sapthesis per usare "Capitolo" al posto di "Chapter")
\usepackage{hyperref}				% Rende cliccabili indice e citazioni
\hypersetup{
    colorlinks = true,				% link colorati invece che con riquadro
    linkcolor  = black,			% link dell'indice: neri (standard per tesi)
    citecolor  = blue,				% link delle citazioni: blu
    urlcolor   = blue
}

%SetFonts

%SetFonts


\title{Introduzione di modelli markoviani e tecniche di apprendimento per l'analisi delle fratture pelviche}
\author{La Rocca Mariastella Gioia}
\IDnumber{2115586}
\authoremail{larocca.2115586@studenti.uniroma1.it}
\AcademicYear{2025/2026}
\copyyear{2026}


\course{Scienze Matematiche per l'Intelligenza Artificiale}      
\courseorganizer{Facoltà di Scienze Matematiche, Fisiche e Naturali. Dipartimento di Matematica}    
\advisor{Prof. Marco Isopi}               

%\date{}							% Activate to display a given date or no date

\begin{document}
\frontmatter
\maketitle
\tableofcontents
\mainmatter
\chapter{Introduzione Generale}
\section{Inquadramento clinico e Gestione del Trauma Pelvico}
Il presente lavoro di tesi si pone l'obiettivo di analizzare l'impatto clinico e prognostico delle fratture pelviche nel contesto dell'emergenza-urgenza. Partendo dalla criticità che tali traumi rappresentano nel setting del Pronto Soccorso, lo studio propone un approccio metodologico innovativo che integra tecniche di apprendimento automatico (Machine Learning) e modellazione stocastica markoviana.

Attraverso un'analisi empirica condotta su un dataset di 991 pazienti, la tesi mira a:
\begin{enumerate}
	\item Identificare i predittori clinici più significativi (come ISS, Glasgow Coma Scale ed Emoglobina) mantenendo la massima interpretabilità medica dei parametri.

	\item Superare la staticità degli score tradizionali (ISS, RTS) attraverso la costruzione di un modello di Markov, capace di descrivere in termini probabilistici l’evoluzione dinamica del paziente e le sue transizioni tra diversi stati di gravità.
\end{enumerate}
L'integrazione di questi modelli si propone come uno strumento di supporto alle decisioni cliniche, offrendo una visione non solo istantanea, ma evolutiva della traiettoria del paziente traumatizzato.

Le fratture pelviche, sebbene statisticamente meno frequenti rispetto ad altre lesioni traumatiche, comportano un elevato carico economico \cite{demetriades2002pelvic} e sono associate a tassi significativi di morbilità e mortalità. Il decesso è spesso determinato dalla gravità delle lesioni associate e dai meccanismi traumatici ad alta energia cinetica che caratterizzano tali eventi.\\

Nonostante i predittori di mortalità siano stati ampiamente documentati in letteratura, la loro applicabilità nel contesto specifico dell'istituto oggetto di studio non è stata aggiornata con dati locali recenti (l’ultimo studio regionale di riferimento risale al periodo 2001-2008). Poiché i protocolli di gestione del trauma (Trauma Management) si sono evoluti drasticamente nell'ultimo decennio, si rende necessario un aggiornamento della capacità prognostica di tali marker.\\

Il presente studio si propone quindi di sviluppare e validare internamente un modello predittivo — in conformità con le linee guida TRIPOD \cite{collins2015tripod} — per aggiornare il valore prognostico dei predittori identificati (tramite la procedura di selezione LASSO) all'interno del contesto istituzionale corrente. L'integrazione successiva di un modello stocastico markoviano permetterà di estendere questa capacità predittiva oltre la dimensione statica, analizzando la dinamicità delle transizioni cliniche del paziente nel tempo.\\

Le fratture del bacino rappresentano una sfida clinica complessa a causa della varietà anatomica e della gravità delle complicanze associate. Tali lesioni possono coinvolgere diverse strutture dell'anello pelvico, tra cui la sinfisi pubica, l’osso innominato, l’acetabolo, l’articolazione sacro-iliaca o il sacro.

Dal punto di vista biomeccanico e clinico, si osserva uno spettro di gravità estremamente ampio:
\begin{itemize}
	\item Lesioni stabili: spesso derivanti da cadute a bassa energia (frequenti nella popolazione geriatrica presente nel dataset), che richiedono solitamente un trattamento sintomatico.

	\item Lesioni instabili: causate da traumi ad alta energia cinetica, caratterizzate da dislocazioni drammatiche che possono compromettere l'integrità vascolare.
\end{itemize}

La principale criticità di queste fratture non risiede solo nel danno osseo, ma nel rischio di emorragie massive e lesioni associate a carico degli apparati genito-urinario, gastrointestinale e del sistema nervoso periferico. La tempestività diagnostica è garantita dall'impiego della radiografia (RX) e, in modo più definitivo, dalla tomografia computerizzata (TC), che costituisce la base per il calcolo dei punteggi di gravità anatomica (ISS). \\

La gestione terapeutica varia in base alla stabilità emodinamica del paziente: mentre le fratture lievi seguono percorsi conservativi, i casi di instabilità con emorragia importante richiedono interventi d'urgenza come la fissazione esterna, l'embolizzazione angiografica o la riduzione aperta con fissazione interna (ORIF). Questa variabilità nei decorsi clinici e la possibilità di repentini peggioramenti emodinamici rendono il trauma pelvico un ambito d'elezione per l'applicazione di modelli stocastici, capaci di mappare le probabilità di transizione tra stati di stabilità e criticità. \\


Come evidenziato dai protocolli internazionali di gestione del trauma (ATLS) \cite{atls2018}, la criticità delle fratture pelviche risiede principalmente nel rischio di shock emorragico massivo. L'emorragia retroperitoneale, spesso associata a lesioni dei plessi venosi o delle arterie iliache, rappresenta la prima causa di mortalità evitabile nelle prime ore dal trauma.

La gestione del paziente segue un percorso decisionale dinamico che può essere schematizzato in fasi sequenziali:
\begin{itemize}

	\item Emostasi temporanea: impiego di fasce pelviche (pelvic binders) per ridurre il volume pelvico.

	\item Stabilizzazione emodinamica: somministrazione di fluidi ed emocomponenti.

	\item Procedure invasive: nei casi di instabilità persistente, si ricorre all'angiografia con embolizzazione o al pelvic packing chirurgico.
\end{itemize}

\subsection{Il razionale del modello stocastico}:
Questa sequenza di interventi non è lineare, ma dipende strettamente dalla risposta del paziente alle manovre rianimatorie. Un paziente può 'transitare' da uno stato di stabilità relativa a uno di instabilità emodinamica severa in qualsiasi momento del percorso diagnostico. Pertanto, la valutazione tramite score statici (come ISS o TRISS) risulta insufficiente per catturare questa evoluzione. L'impiego di una Catena di Markov in questo studio permette di modellare proprio queste transizioni critiche, fornendo una stima probabilistica del rischio di peggioramento basata sulla variazione in tempo reale di parametri quali la Pressione Sistolica e i Lattati ematici, indicatori indiretti ma precoci dello stato di emorragia occulta.

\section{Indici Descrittivi e Sistemi di Valutazione Tradizionali}
\subsection{Indici Descrittivi}

La valutazione della gravità del paziente con frattura pelvica si avvale di una serie di indici descrittivi che possono essere suddivisi in parametri fisiologici, biochimici e anatomici. Questi indici non sono solo descrittivi, ma fungono da marker prognostici precoci.

\subsubsection{Parametri Fisiologici e Vitali}
I parametri vitali rappresentano la risposta immediata dell'organismo al trauma e al relativo insulto emorragico:
\begin{itemize}
	\item {Pressione Arteriosa Sistolica (Psis)}: È l'indicatore principale della stabilità emodinamica. Valori inferiori a 90 mmHg definiscono convenzionalmente lo stato di shock e sono associati a una maggiore probabilità di transizione verso stati critici nel modello stocastico.
	\item {Frequenza Cardiaca (FC)}: Segnala il tentativo di compenso del sistema cardiocircolatorio. La tachicardia è spesso il primo segno di ipovolemia, precedendo la caduta della pressione.
	\item {Glasgow Coma Scale (GCS)}: Valuta lo stato di coscienza e l'integrità del sistema nervoso centrale, parametro fondamentale per escludere lesioni cerebrali associate che complicano il quadro clinico.
\end{itemize}

\subsubsection{Marker Biochimici e Metabolici}
A differenza dei parametri vitali, i marker biochimici riflettono lo stato di perfusione dei tessuti:
\begin{itemize}
\item {Lattati (Lac)}: Rappresentano il prodotto del metabolismo anaerobico cellulare. Come emerso dall'analisi descrittiva di questo studio, i Lattati costituiscono uno dei predittori più sensibili di mortalità \cite{mcnelis2001lactate}. Un loro incremento riflette un'ipossia tissutale persistente, rendendoli un parametro dinamico ideale per monitorare l'evoluzione del paziente.
\item {Eccesso di Basi (Base Excess - BE)}: Misura l'entità dello squilibrio acido-base, fornendo una stima indiretta della gravità dell'emorragia e della necessità di trasfusioni massive.
\end{itemize}

\paragraph{Indici Anatomici di Gravità}
L'indice principe è l'Injury Severity Score (ISS) \cite{baker1974iss}, che quantifica il carico lesivo globale. Nel contesto delle fratture pelviche, l'ISS permette di pesare non solo la frattura ossea ma anche le lesioni d'organo associate, fornendo una misura della complessità del trauma che rimane costante durante tutto l'iter ospedaliero.

\subsection{Sistemi di Valutazione Tradizionali}
Per standardizzare la valutazione della gravità e permettere un confronto tra diverse casistiche cliniche, la medicina d'emergenza si avvale di score validati internazionalmente. I tre pilastri della valutazione tradizionale sono:

\begin{enumerate}
\item {Injury Severity Score (ISS)}: È lo score di riferimento per la gravità {anatomica}. Viene calcolato assegnando un punteggio AIS (Abbreviated Injury Scale) a sei regioni corporee e sommando i quadrati dei tre punteggi più alti. Nel trauma pelvico, l'ISS è fondamentale per quantificare l'entità del danno strutturale, ma presenta il limite di essere un parametro statico, calcolabile solo dopo completamento diagnostico radiologico.

\item {Revised Trauma Score (RTS)} \cite{champion1989rts}: È uno score {fisiologico} basato su tre parametri rilevati al triage: Glasgow Coma Scale (GCS), Pressione Arteriosa Sistolica (Psis) e Frequenza Respiratoria. L'RTS è un indicatore immediato della stabilità del paziente, ma può variare rapidamente durante la permanenza in Pronto Soccorso, rendendo necessaria una modellazione dinamica delle sue variazioni.

\item {Trauma and Injury Severity Score (TRISS)} \cite{boyd1987triss}: Rappresenta il tentativo di integrare le informazioni anatomiche (ISS) e fisiologiche (RTS) con l'età del paziente. Attraverso una formula di regressione logistica, il TRISS fornisce una stima della probabilità di sopravvivenza ($P_s$). Sebbene sia stato per decenni il riferimento globale, il TRISS non tiene conto di marker biochimici moderni come i lattati, né della natura stocastica delle transizioni cliniche.
\end{enumerate}

L'analisi condotta in questa tesi utilizzerà questi score come base di partenza per identificare, tramite tecniche di apprendimento automatico, i predittori più significativi all'interno del dataset specifico in esame.

\section{Obiettivo della Tesi: un approccio statistico-predittivo}
L’obiettivo centrale di questo lavoro è lo sviluppo di un modello metodologico che integri l’analisi empirica dei dati clinici con tecniche avanzate di modellazione stocastica, al fine di migliorare la capacità predittiva e la comprensione dinamica del trauma pelvico in fase acuta.\\

L'analisi predittiva consiste nell’utilizzare dati, algoritmi statistici e tecniche di machine learning per individuare la probabilità di risultati futuri basandosi sui dati storici. L'obiettivo è andare oltre la comprensione di "cosa è successo" per arrivare a una valutazione proattiva di "quello che accadrà". In ambito traumatologico, ciò si traduce nella capacità di anticipare l'evoluzione del paziente partendo dai segnali deboli rilevati al triage.

I modelli predittivi sfruttano risultati noti per sviluppare (o addestrare) un modello applicabile a scenari nuovi. La modellizzazione fornisce risultati sotto forma di previsioni che rappresentano la probabilità di una variabile target (nel nostro caso, l'esito clinico o la transizione tra stati di gravità) in base alla significatività stimata a partire da un insieme di variabili di input.\\

Esistono due macro-tipologie di modelli predittivi:
\begin{itemize}
	\item Modelli di classificazione: prevedono l'appartenenza a una classe discreta. In questo studio, la classificazione principale riguarda l'esito finale (0 per il sopravvissuto, 1 per il decesso) o la categorizzazione del rischio (basso, medio, alto).

	\item Modelli di regressione: prevedono un valore numerico continuo. Sebbene la regressione lineare sia destinata a dati che seguono una distribuzione normale, in ambito clinico la regressione logistica risulta spesso più appropriata, poiché permette di gestire variabili di risposta categoriche (esito binario) in base a predittori misti \cite{james2013introduction}.
\end{itemize}

\subsection{Sviluppo futuro Catena di Markov}
Oltre alla distinzione tra modelli di classificazione e regressione, questo studio introduce l'impiego delle Catene di Markov come strumento per superare i limiti della previsione puntuale e statica. Mentre la classificazione logistica permette di categorizzare l'esito finale del paziente (Sopravvissuto/Deceduto), la modellazione markoviana consente di mappare l'evoluzione dinamica del trauma lungo l'asse temporale della gestione in emergenza.

Tale approccio si fonda sulla cosiddetta proprietà di Markov, secondo la quale la probabilità di transizione del sistema verso uno stato futuro dipende esclusivamente dallo stato clinico attuale del paziente, indipendentemente dalla sequenza di eventi che hanno condotto a tale condizione. Matematicamente, questo si esprime come:\\

$P(X_{n+1}=j \mid X_n=i, X_{n-1}=i_{n-1}, \dots, X_0=i_0) = P(X_{n+1}=j \mid X_n=i)$\\

L'integrazione di questo concetto crea un ponte teorico fondamentale tra la statistica descrittiva e l'analisi predittiva: l'esito clinico non viene più interpretato come un evento isolato, ma come il risultato finale di una traiettoria stocastica in cui il paziente transita tra diversi livelli di gravità (stati). In questo modo, la tesi evolve da una semplice analisi dei fattori di rischio a una vera e propria simulazione del decorso clinico del traumatizzato.\\

Il cuore dell'approccio proposto in questa tesi si articola in tre fasi sequenziali. Si parte con un'ottimizzazione della qualità del dato (Data Cleaning), più precisamente superare l'ambiguità delle registrazioni cliniche attraverso una distinzione rigorosa tra valori nulli e dati mancanti. Questo passaggio è propedeutico a qualsiasi inferenza statistica affidabile, garantendo che il modello non venga falsato da "zeri" che rappresentano in realtà la mancata rilevazione di un parametro (come emerso per i Lattati o la Pressione Sistolica).\

Successivamente avviene una selezione dei predittori tramite Machine Learning (LASSO): Identificare, all'interno della vasta mole di variabili rilevate (36 parametri), il set minimo di predittori dotati del massimo potere informativo. Si è scelto di privilegiare la tecnica LASSO (Least Absolute Shrinkage and Selection Operator) \cite{tibshirani1996lasso} rispetto alla PCA. Questo permette di isolare parametri critici — come l'ISS e i Lattati — mantenendo intatta la loro interpretabilità clinica e la loro unità di misura originale, rendendo i risultati immediatamente decodificabili in un contesto di emergenza.\\

Infine si tratta una modellazione stocastica del decorso clinico (Catene di Markov): Rappresenta il traguardo innovativo dello studio. L'obiettivo è superare la staticità degli score tradizionali per mappare la "traiettoria" del paziente. Attraverso la costruzione di una Catena di Markov, lo studio mira a calcolare le probabilità di transizione tra diversi stati di gravità clinica, definiti sulla base dei predittori selezionati. Questo approccio permette di quantificare il rischio di peggioramento emodinamico in funzione del tempo e dello stato attuale del paziente.

In sintesi, la tesi si propone di trasformare un dataset statico di Pronto Soccorso in uno strumento predittivo dinamico. L'integrazione tra la precisione selettiva della LASSO e la natura evolutiva dei modelli Markoviani punta a fornire una base matematica solida per lo sviluppo di futuri sistemi di supporto alle decisioni cliniche (Clinical Decision Support Systems), capaci di ottimizzare i tempi di intervento e la gestione delle risorse nella cura del trauma pelvico.



\chapter{Descrizione del Campione e Analisi Esplorativa dei Dati (EDA)}
\section{Profilo del Campione e Provenienza dei Dati}
La presente analisi è stata condotta su un dataset retrospettivo comprendente 991 pazienti giunti in emergenza per trauma pelvico. I dati riflettono la realtà clinica di un reparto di Pronto Soccorso di II livello.

La distribuzione di genere mostra una lieve prevalenza femminile, con 550 donne (56,1\%) e 431 uomini (43,9\%). 

L’età media del campione è di 58,45 anni (DS ± 22,72), con un intervallo compreso tra 15 e 102 anni. Tale dato è di estremo interesse clinico: la deviazione standard elevata riflette la natura "bimodale" del trauma pelvico. Come evidenziato dall'analisi dei dati, il campione include sia pazienti giovani vittime di traumi ad alta energia, sia una consistente quota di pazienti geriatrici soggetti a fratture da fragilità ossea (osteoporosi).

\subsection{Meccanismo del Trauma e Priorità Clinica}

L’analisi dei meccanismi lesivi (Tabella 2.1) conferma l'eterogeneità del campione. I due scenari predominanti risultano essere:
\begin{itemize}
\item {Incidente stradale (Codice 2)}: con 382 casi, rappresenta la causa principale di accesso, tipicamente associata a dinamiche ad alta energia cinetica.
\item {Incidente domestico (Codice 4)}: con 234 casi, rappresenta la seconda causa, coerente con l'età media elevata rilevata e indicativa di cadute a bassa energia in ambiente protetto.
\end{itemize}

Questi due meccanismi coprono complessivamente il 62,1\% della casistica. In termini di gravità percepita al momento dell'arrivo, il 60,6\% dei pazienti è stato classificato con Triage Rosso (Codice 1). Questo dato sottolinea l'elevata criticità del trauma pelvico nel setting del Pronto Soccorso, dove la stabilità emodinamica è spesso compromessa fin dall'ingresso.

\begin{table}[h]
\centering
\caption{Caratteristiche demografiche e cliniche del campione (N=991)}
\begin{tabular}{ll}
\hline
\textbf{Variabile} & \textbf{Valore (Media $\pm$ DS o \%)} \\ \hline
Età (anni) & 58,45 $\pm$ 22,72 \\
Genere Femminile & 550 (56,1\%) \\
Genere Maschile & 431 (43,9\%) \\
Triage Rosso (Codice 1) & 601 (60,6\%) \\
Meccanismo: Incidente Stradale & 382 (38,5\%) \\
Meccanismo: Caduta Domestica & 234 (23,6\%) \\ \hline
\end{tabular}
\end{table}

Un elemento cruciale nella descrizione del campione è la distribuzione dell'esito finale. Nel dataset analizzato, il tasso di sopravvivenza è del 87\%, mentre il tasso di mortalità è del 12\%. Questa proporzione definisce il bilanciamento del dataset: una mortalità significativa (tipica del trauma pelvico complesso) permette al modello di apprendimento di identificare con maggiore precisione le caratteristiche dei pazienti ad alto rischio.

Oltre ai dati demografici, il campione si caratterizza per un'elevata gravità lesiva globale: l'Injury Severity Score (ISS) medio di 23,2 conferma che la maggior parte dei pazienti appartiene alla categoria del 'trauma maggiore' (ISS $> 15$). In particolare, una quota rilevante di pazienti presenta un ISS superiore a 34, valore che correla con un rischio di decesso estremamente elevato. Il valore di Glasgow Coma Scale (GCS) medio riflette lo stato neurologico all'ingresso, con una quota di pazienti che presenta punteggi inferiori a 8, indicativi di coma o grave compromissione cerebrale associata al trauma.

L'instabilità del campione è ulteriormente testimoniata dalla necessità di manovre rianimatorie invasive immediate: l'11\% dei pazienti (n = 110) ha richiesto l'intubazione orotracheale (IOT) direttamente in Pronto Soccorso o nelle primissime fasi della gestione. Questo parametro è un indicatore indiretto di insufficienza respiratoria o shock profondo, e rappresenta uno degli stati critici che verranno analizzati nel modello stocastico.

\section{Data Cleaning: La Gestione del "Dato Mancante"}
È un assioma ormai consolidato nel campo della Data Science e dell'apprendimento automatico che la maggior parte dell'impegno necessario per raggiungere un obiettivo analitico risieda nella fase di pulizia dei dati. Spesso si stima, con una punta di ironia ma con grande realismo, che tale attività assorba ben oltre il 50\% del tempo totale di ricerca. Nella pratica clinica, questa percentuale è spesso persino superiore, data la complessità e l'eterogeneità delle registrazioni in regime di emergenza.

In questo studio, si è partiti dalla consapevolezza che ogni dataset reale è intrinsecamente "sporco" e che l'aspirazione a un database perfettamente pulito è spesso un'utopia. L'obiettivo metodologico non è stato quindi il raggiungimento di una purezza teorica, quanto piuttosto la riduzione del rumore e delle incoerenze per trasformare il dato grezzo in un dato "meno sporco" e, soprattutto, utilizzabile. 
L'obiettivo della metodologia qui descritta è stato trasformare questo dataset grezzo in una base analitica dotata di utilità clinica.

A tal proposito, è fondamentale distinguere tra qualità del dato (accuratezza formale) e utilità del dato (capacità informativa). Un dataset può presentare criticità qualitative — come l'elevata frequenza di dati mancanti riscontrata nei Lattati in questa casistica — ma mantenere un'enorme utilità prognostica. Il processo di Data Cleaning descritto di seguito è stato quindi guidato dalla necessità di preservare l'utilità clinica dei parametri, correggendo al contempo le anomalie sistemiche che avrebbero potuto introdurre bias nei modelli predittivi e markoviani.\\

La qualità dei risultati di un modello predittivo, in particolare quando si utilizzano approcci stocastici come le Catene di Markov, dipende in modo critico dalla pulizia e dall'integrità del database di partenza. Come sottolineato da Ganti e Das Sarma \cite{ganti2013datacleaning}, il Data Cleaning non deve essere inteso come una semplice correzione di errori isolati, ma come un processo strutturato volto a garantire che il dato rifletta fedelmente la realtà clinica, evitando che anomalie tecniche si trasformino in bias decisionali.

In questa tesi, il processo di preparazione dei dati è stato suddiviso in tre fasi principali: Data Profiling, Data Standardization e Outlier Treatment.

\subsection{Data Profiling e Identificazione dei Dati Mancanti}
La prima fase ha riguardato il "profiling" del dataset, ovvero l'analisi statistica preliminare per valutare la qualità delle informazioni raccolte. Una delle sfide maggiori è stata la gestione del valore "0". In un contesto di Pronto Soccorso, il valore zero può essere ambiguo: può rappresentare un dato clinico reale (es. assenza di una comorbidità) o, più frequentemente, l'omissione di una registrazione (dato mancante).\\

Ad esempio, analizzando la colonna della Pressione Arteriosa Sistolica (Psis), sono stati riscontrati inizialmente 74 valori pari a zero. Matematicamente, una pressione pari a zero in un paziente giunto al triage indicherebbe un decesso istantaneo, rendendo inutile ogni altra misurazione. Poiché tali record presentavano invece altri parametri vitali (come la frequenza cardiaca), si è dedotto per inferenza logica che lo "0" non fosse un valore misurato, ma un dato non registrato.

Seguendo le raccomandazioni sulla standardizzazione del dato, è stata operata una distinzione rigorosa: per tutti i parametri fisiologici (Pressione Arteriosa, Frequenza Cardiaca, GCS) e biochimici (Lattati, Emoglobina), si è proceduto alla conversione sistematica degli zeri in NaN (Not a Number). Questa scelta è fondamentale per non falsare le medie cliniche: includere un valore di pressione sistolica pari a 0 nel calcolo della stabilità emodinamica avrebbe portato a una sovrastima sistematica della gravità e a un’errata calibrazione delle probabilità di transizione nel modello di Markov.\\

\textbf{Nota sui conflitti di formato e Outlier:} 
Un'altra criticità pratica ha riguardato la variabile Età. Durante il \textit{data profiling}, è emerso che alcuni record erano stati interpretati dal software di database come date anziché come numeri interi. Questo errore sistematico avrebbe reso impossibile qualsiasi calcolo statistico. Si è reso necessario un intervento di normalizzazione del formato, coerentemente con i compiti di \textit{Data Standardization}, forzando la conversione in valori numerici e recuperando l'integrità della distribuzione demografica, che ha poi rivelato un'età media di 58 anni.\\
 
Parallelamente, per la variabile Lattati (Lac), è stato identificato e rimosso un valore outlier di 765 mmol/L, palesemente frutto di un errore di digitazione in quanto clinicamente incompatibile con la vita. A seguito di questa rimozione, la media dei Lattati è passata da 10.2 mmol/L a un valore più rappresentativo della realtà clinica del campione. L'analisi dei dati mancanti rivela che il parametro dei Lattati è presente solo nel 23\% dei casi (n=225). Nonostante l'elevata frammentarietà, si è deciso di mantenere la variabile per la sua nota rilevanza prognostica, procedendo con analisi specifiche sui casi disponibili per la fase di \textit{Feature Selection}.

L'analisi esplorativa ha infatti mostrato che nei pazienti deceduti la media dei lattati era significativamente più alta (5.52 vs 2.39 mmol/L). Questo suggerisce che, sebbene il dato sia scarso (bassa qualità quantitativa), esso è estremamente prezioso (alta utilità qualitativa) perché viene raccolto dai medici quasi esclusivamente nei casi di sospetta instabilità emodinamica. Scartare i lattati avrebbe significato privare il modello di Markov di uno dei suoi predittori più sensibili.


\subsection{Codifica delle Variabili (Feature Engineering)}

Una volta garantita l'integrità qualitativa dei dati attraverso il cleaning, è stato necessario procedere alla loro trasformazione semantica per renderli compatibili con gli algoritmi di Machine Learning e con la successiva modellizzazione markoviana. Questo processo, noto come Feature Engineering, permette di convertire informazioni categoriali in formati numerici processabili matematicamente senza perdere il contenuto informativo originale.

Nello specifico, sono state adottate due strategie di codifica:
\begin{itemize}
\item Label Encoding: utilizzata per le variabili binarie (come \textit{Sesso} e \textit{IOT}) e per le scale ordinali. Questa tecnica assegna un valore numerico (0 o 1) a ciascuna categoria. Ad esempio, per la variabile IOT, il valore 0 indica l'assenza di intubazione e 1 la sua esecuzione, permettendo alla regressione LASSO di pesare l'impatto di tale intervento sulla sopravvivenza.

\item One-Hot Encoding: applicata alle variabili categoriali non ordinate, con particolare riferimento al \textit{Meccanismo del trauma}. Poiché non esiste una gerarchia intrinseca tra un "incidente stradale" e una "precipitazione", l'uso di una semplice scala numerica (1, 2, 3...) avrebbe indotto l'algoritmo a ipotizzare erroneamente un ordine di grandezza o di importanza. La tecnica One-Hot Encoding supera questo limite creando variabili "dummy" separate per ogni meccanismo, garantendo che ogni dinamica traumatica venga valutata in modo indipendente e oggettivo.
\end{itemize}

Questa fase di trasformazione rappresenta l'ultimo anello della catena di preparazione del dato, fungendo da ponte tra l'analisi esplorativa descrittiva e la fase di modellazione predittiva

\section{Analisi Descrittiva Univariata}
Dopo la fase di pulizia e codifica, è stata condotta un'analisi statistica per esplorare la natura dei parametri clinici e la loro relazione con l'esito. Questa fase permette di identificare i pattern primari e di convalidare la scelta dei predittori.

L'analisi della distribuzione dell'età (Figura 2.1) conferma la natura eterogenea del campione. L'istogramma evidenzia una distribuzione "bimodale": un primo picco nella fascia giovane-adulta (20-40 anni), prevalentemente correlato a incidenti stradali e traumi ad alta energia; un secondo picco più pronunciato nella fascia geriatrica (>70 anni), legato a fratture da fragilità per cadute domestiche.

\begin{figure}[htbp]
    \centering
    \includegraphics[width=0.8\textwidth]{figura_2_1.png}
    \caption{L’istogramma illustra la distribuzione anagrafica dei 991 pazienti, evidenziando la natura bimodale del trauma pelvico. Si osservano due picchi distinti: il primo nella fascia giovane-adulta (20-40 anni), solitamente legata a incidenti stradali ad alta energia, e il secondo nella fascia geriatrica (sopra i 70 anni), riconducibile a fratture da fragilità per cadute domestiche. Questa distribuzione conferma la necessità di considerare l’età come un fattore determinante nell'analisi del rischio.}
    \label{fig:mio_grafico}
\end{figure}



Questa distribuzione giustifica l'inserimento dell'età come variabile nel modello, poiché la riserva fisiologica del paziente varia drasticamente tra queste due coorti.

Lo studio della distribuzione delle variabili chiave mette in luce la natura complessa del campione analizzato:
\begin{itemize}

	\item {ISS (Injury Severity Score)}: L'ISS si conferma come il principale indicatore di gravità. La distribuzione globale (Figura 2.2) mostra una media di 23.2 e una mediana di 22.
	
	\begin{figure}[htbp]
    \centering
    \includegraphics[width=0.8\textwidth]{figura_2_2.png}
    \caption{Il grafico mostra la frequenza dei punteggi di gravità anatomica (ISS) all'interno della coorte. La distribuzione appare sbilanciata verso destra, con una media di 23.2 (linea rossa tratteggiata) e una mediana di 22 (linea nera). Poiché oltre il 60\% del campione presenta un ISS superiore a 15, il dato conferma che la casistica del Pronto Soccorso in esame è composta prevalentemente da 'traumi maggiori', caratterizzati da lesioni multi-organo complesse oltre alla frattura pelvica}
    \label{fig:mio_grafico}
\end{figure}

Nel confronto bivariato (Boxplot in Figura 2.3), emerge una separazione netta: i pazienti sopravvissuti presentano una mediana di ISS pari a 18 (trauma moderato.grave); i pazienti deceduti mostrano una mediana di 34, con gran parte del gruppo concentrata sopra la soglia critica di 25. 

Questo "gap" visivo tra le scatole del boxplot è la prova statistica del forte potere predittivo dell'ISS per la mortalità nel trauma pelvico.

\begin{figure}[htbp]
    \centering
    \includegraphics[width=0.8\textwidth]{figura_2_3.png}
    \caption{Il confronto bivariato evidenzia la distribuzione dell'ISS scorporata per esito clinico. Si osserva un netto slittamento della popolazione dei deceduti (area rossa) verso i valori alti della scala, con una mediana di 34, rispetto ai sopravvissuti (area blu) che presentano una concentrazione sui traumi moderati (mediana 18). La scarsa sovrapposizione delle due distribuzioni conferma visivamente l'ISS come il principale fattore discriminante della mortalità.}
    \label{fig:mio_grafico}
\end{figure}

	\item {Lattati (Lac)}: L'analisi dei Lattati, condotta sui 225 pazienti in cui il dato era disponibile, rivela informazioni cruciali. Nonostante la rimozione dell'outlier estremo (765 mmol/L), la distribuzione rimane fortemente asimmetrica verso destra.
La mediana dei sopravvissuti è di 2.39 mmol/L, vicina ai range di normalità. Al contrario, il gruppo dei decessi presenta una mediana di 5.52 mmol/L, valore che indica uno stato di acidosi lattica e ipoperfusione tissutale severa. Tale differenza (più che raddoppiata) indica che il metabolismo anaerobico è un segnale precoce di esito infausto, spesso precedendo la caduta della pressione arteriosa.

	\item {Pressione Sistolica (Psis)}: La Pressione Arteriosa Sistolica all'ingresso rappresenta il parametro clinico di riferimento per la valutazione immediata dello stato emodinamico. Nel campione analizzato, il valore medio si attesta a {122.13 mmHg}. Sebbene tale dato indichi una stabilità media della coorte, la deviazione standard riflette la presenza di un sottogruppo di pazienti in stato di ipotensione severa (Psis < 90 mmHg). Come evidenziato nell'analisi bivariata, questo parametro subisce un calo significativo nei pazienti con esito infausto, confermandosi come un indicatore precoce e fondamentale di shock emorragico scompensato.
\end{itemize}


\section{Analisi Bivariata: Correlazione con l'Esito}

\begin{figure}[htbp]
    \centering
    \includegraphics[width=0.8\textwidth]{figura_2_4.png}
    \caption{Matrice di correlazione di Spearman. Il grafico visualizza il grado di associazione tra i parametri clinici e l'esito. Si notano le forti correlazioni tra i marker metabolici e la gravità anatomica (ISS), fornendo la base per la successiva selezione dei predittori tramite tecnica LASSO.}
    \label{fig:mio_grafico}
\end{figure}

Al fine di identificare le variabili con il maggior potere discriminante rispetto all'esito clinico, è stata condotta un'analisi comparativa tra il gruppo dei sopravvissuti e quello dei deceduti.

Dall'analisi emerge che i parametri con la maggiore significatività statistica sono l'ISS e i Lattati. In particolare, i pazienti con esito infausto presentano valori di ISS mediani sensibilmente più elevati, confermando come la gravità anatomica del trauma pelvico sia un predittore primario di mortalità. Parallelamente, il valore dei Lattati all'ingresso si conferma un marker biochimico cruciale: valori elevati riflettono uno stato di ipoperfusione tissutale che correla direttamente con la probabilità di decesso.

Per visualizzare l'interdipendenza tra i parametri senza ricorrere alla riduzione dimensionale (PCA), è stata elaborata una matrice di correlazione di Spearman (Figura 2.4). 



Dalla Heatmap emergono tre cluster di variabili:
\begin{enumerate}
\item Cluster Emodinamico: Una forte correlazione negativa tra Lattati e Pressione Sistolica ($r = - 0.42$); al diminuire della pressione, i lattati aumentano, descrivendo matematicamente lo shock emorragico.
\item Cluster Metabolico: Una correlazione strettissima tra Lattati, pH e Bicarbonati (Hco3-), che conferma la coerenza interna dei dati biochimici raccolti.
\item Cluster Prognostico: L'Esito clinico mostra le correlazioni più forti con l'ISS, i Lattati e il GCS, confermando questi parametri come i "predittori significativi".
\end{enumerate}


\chapter{Selezione dei Predittori (Feature Selection)}

\section{Obiettivo della selezione}
L’elevata dimensionalità del dataset originale (36 variabili per 991 pazienti) richiede una riduzione del numero di predittori per evitare l'overfitting e garantire la robustezza del modello finale. L'obiettivo è isolare esclusivamente i parametri dotati di reale valore prognostico, eliminando il "rumore" di fondo prodotto da variabili poco influenti o ridondanti.

Si è evitato l'utilizzo dell'Analisi delle Componenti Principali (PCA). Sebbene la PCA sia efficace nel ridurre la dimensionalità, essa genera nuove variabili (componenti) che sono combinazioni lineari dei valori misurati, rendendo impossibile l'interpretazione clinica diretta. Per questa tesi è invece fondamentale mantenere l'identità fisica e medica di ogni parametro (es. mmHg per la pressione, mmol/L per i lattati) in vista della costruzione del modello stocastico e per poter definire in modo trasparente gli stati della Catena di Markov.

\subsubsection{Il ruolo della Feature Selection nei modelli predittivi}
La selezione dei predittori (Feature Selection) rappresenta un pilastro fondamentale nel processo di costruzione di un modello di apprendimento automatico, specialmente in ambito clinico. L'obiettivo non è semplicemente ridurre il numero di variabili, ma ottimizzare l'architettura informativa del modello per quattro ragioni essenziali:

\begin{enumerate}
\item {Miglioramento dell'accuratezza (Improved Accuracy)}: L'inclusione di variabili ridondanti o irrilevanti agisce come "rumore" statistico. Concentrando l'addestramento dell'algoritmo esclusivamente sulle feature dotate di reale potere segnaletico, il modello è in grado di apprendere pattern più robusti, riducendo il rischio di {overfitting} e migliorando la capacità di generalizzazione su nuovi pazienti.
\item {Efficienza computazionale (Faster Training)}: Sebbene la potenza di calcolo attuale permetta di gestire dataset di medie dimensioni, la riduzione dello spazio delle feature accelera significativamente i tempi di elaborazione. Questo aspetto diventa cruciale nella prospettiva di un'implementazione in tempo reale del modello in contesti di emergenza, dove la velocità di risposta è un requisito operativo.

\item {Massima interpretabilità (Greater Interpretability)}: In medicina la "black box" non è accettabile. Un modello con pochi input significativi permette al clinico di comprendere la logica sottostante alla previsione. Sapere che il rischio è guidato da ISS e Lattati consente un'azione mirata, a differenza di modelli complessi dove l'impatto dei singoli parametri fisici viene oscurato.

\item {Superamento della "Maledizione della Dimensionalità" (Avoiding the Curse of Dimensionality)}: All'aumentare del numero di variabili, lo spazio dei dati diventa estremamente sparso. Con un campione di 991 pazienti, l'utilizzo di tutte le 36 variabili originali rischierebbe di rendere il modello instabile. La Feature Selection riduce questa complessità, garantendo che il numero di osservazioni sia statisticamente sufficiente rispetto ai parametri stimati.

\end{enumerate}
In questo contesto, la regressione Lasso si inserisce come un metodo di Embedded Feature Selection, integrando la selezione dei predittori direttamente all'interno del processo di stima del modello.

\section{Il problema della dimensionalità: perché evitare la PCA}
La mole di dati raccolta in regime di Pronto Soccorso, sebbene estremamente preziosa, pone una sfida metodologica nota come "maledizione della dimensionalità" (curse of dimensionality). Nel dataset in esame, a fronte di una popolazione di 991 pazienti, sono state registrate inizialmente 36 variabili potenziali. In statistica, un numero eccessivo di predittori rispetto alle osservazioni può portare a modelli instabili, difficili da generalizzare e inclini all'overfitting.

\subsection{I limiti dell'Analisi delle Componenti Principale (PCA)}
Una delle tecniche più diffuse per ridurre la dimensionalità è l'Analisi delle Componenti Principali (PCA). La PCA trasforma l'insieme originale di variabili correlate in un set ridotto di variabili incorrelate, chiamate componenti principali, attraverso una procedura di ortogonalizzazione.

Nonostante l'efficacia matematica, l'applicazione della PCA in ambito medico-traumatologico presenta limiti critici:
\begin{itemize}
\item {Perdita di identità clinica}: La PCA genera nuove variabili che sono combinazioni lineari dei valori misurati. Una "Componente 1" composta da una frazione di Pressione Sistolica, una di Età e una di Lattati non possiede un'unità di misura fisica né un significato clinico immediato.
\item {Difficoltà di applicazione al Triage}: In un contesto di emergenza, il medico ha bisogno di agire su parametri reali (es. "il paziente è ipoteso" o "i lattati sono alti"). Un modello basato su componenti astratte risulterebbe inutilizzabile come strumento di supporto decisionale in tempo reale.
\item {Natura mista dei dati}: Come osservato nel dataset, molte variabili sono categoriali o ordinali. La PCA standard è progettata per variabili numeriche continue e non si presta in modo ottimale alla natura mista della casistica clinica del trauma pelvico.
\end{itemize}

Per queste ragioni si è scelto di scartare le tecniche di riduzione della dimensionalità basate su combinazioni lineari. L'obiettivo è stato spostato verso metodi di Feature Selection capaci di mantenere intatta l'integrità delle variabili originali, identificando i predittori "fisici" più significativi. Tale necessità ha condotto alla scelta dei metodi \textit{embedded} e, nello specifico, della regressione LASSO.

\section{Inquadramento metodologico: i Metodi Embedded}
Esistono diverse strategie per effettuare la Feature Selection, raggruppabili principalmente in tre categorie: metodi filter, wrapper ed embedded. Per gli scopi di questa tesi, la scelta è ricaduta sui metodi embedded, in quanto rappresentano l'approccio più avanzato ed efficiente per la gestione di dataset clinici complessi.

I metodi embedded eseguono la selezione dei predittori direttamente durante il processo di addestramento del modello (training). Essi combinano i vantaggi dei metodi filter (velocità e rigore statistico) e dei metodi wrapper (interazione con il modello), integrando la selezione all'interno dell'algoritmo stesso. Questo permette al modello di identificare dinamicamente quali feature siano realmente rilevanti mentre impara a prevedere l'esito.

Tra le tecniche embedded più diffuse nel campo della Data Science si annoverano:
la {Regolarizzazione L1 (LASSO)}, che impone una penalità ai coefficienti, riducendo a zero quelli delle feature meno significative;{Alberi di Decisione e Random Forests}, che selezionano le variabili in base alla riduzione dell'impurità (Gini o Entropia); {Gradient Boosting}, che identifica i predittori capaci di ridurre maggiormente l'errore di previsione ad ogni iterazione.

\subsubsection{Vantaggi e Limitazioni}
L'adozione di un approccio embedded offre significativi vantaggi in termini di efficienza: a differenza dei metodi wrapper, non richiede un carico computazionale eccessivo, pur mantenendo una capacità di apprendimento specifica per il modello utilizzato (Model-specific learning).

Tuttavia, tali metodi presentano anche dei limiti intrinseci. Rispetto ai più semplici metodi filter (come i test statistici univariati), la loro interpretazione può risultare più complessa, richiedendo una profonda comprensione del processo di apprendimento dell'algoritmo. Inoltre, non sono universalmente applicabili, poiché richiedono algoritmi che supportino nativamente tecniche di regolarizzazione o di analisi dell'importanza delle feature.

\section{Metodologia: Regressione LASSO}
Nell'ambito dell'apprendimento automatico (Machine Learning), i modelli di regressione costituiscono strumenti analitici di primaria importanza per la modellazione delle relazioni intercorrenti tra variabili indipendenti (predittori) e variabili dipendenti (esiti). Attraverso lo studio di queste interdipendenze, le tecniche di regressione consentono di generare previsioni accurate su dati non ancora osservati, permettendo di sostituire congetture cliniche soggettive con \textit{insight data-driven} basati sulla significatività statistica stimata dal modello.

La regressione logistica, in particolare, è una tecnica di apprendimento supervisionato che modella la probabilità di appartenenza di un'osservazione a una determinata classe (nel nostro caso, la sopravvivenza o il decesso) \cite{hosmer2013applied}. Tuttavia, quando il numero di variabili di input è elevato, i modelli tradizionali rischiano di incorrere in due criticità sistemiche: l'eccessiva complessità del modello e la ridondanza informativa.

Per ovviare a tali limiti, in questo studio è stata adottata la Regressione LASSO (Least Absolute Shrinkage and Selection Operator) \cite{hastie2009elements}. Questa scelta è stata dettata dalla necessità di bilanciare l'efficienza predittiva con l'esigenza di interpretabilità clinica. A differenza di modelli "ensemble" più opachi come le Random Forest, la LASSO permette di isolare l'impatto di parametri fisici e biochimici (come l'ISS o l'Hb) senza ricorrere a combinazioni lineari astratte, mantenendo intatta l'identità medica di ogni variabile.

\subsection{Meccanismi di Shrinkage e Geometria della LASSO}
Per comprendere appieno l'efficacia della LASSO nella selezione dei predittori, è necessario analizzare il meccanismo di contrazione dei coefficienti (shrinkage). A differenza della regressione OLS, che cerca di minimizzare esclusivamente l'errore sui dati di addestramento, la LASSO impone un 'budget' limitato alla grandezza totale dei coefficienti $\beta$.

\subsubsection{Il processo di selezione}
Il funzionamento può essere sintetizzato in tre passaggi logici:
\begin{enumerate}
\item Penalizzazione: L'algoritmo aggiunge un costo per ogni variabile inclusa nel modello. Se una variabile non contribuisce a ridurre l'errore più di quanto 'costi' in termini di penalità, il suo peso viene ridotto.

\item Sparsità: A causa della natura matematica della norma $L_1$ (basata sul valore assoluto), la contrazione dei coefficienti non è asintotica (come nella regressione Ridge), ma spinge i valori esattamente a zero.

\item Identificazione dei predittori: Al termine del processo, solo le variabili dotate di un forte legame statistico con l'esito mantengono un coefficiente non nullo. Questo trasforma il modello di regressione in un vero e proprio filtro selettivo.

\end{enumerate}

\subsubsection{Interpretazione Geometrica}
Dal punto di vista geometrico, la LASSO può essere visualizzata come l'intersezione tra la funzione di errore e una regione di vincolo a forma di diamante (rappresentata dalla norma $L_1$). Poiché il diamante presenta degli 'spigoli' in corrispondenza degli assi coordinati (dove uno o più coefficienti sono pari a zero), è estremamente probabile che la soluzione ottimale tocchi proprio uno di questi vertici.

È questa caratteristica geometrica a garantire che il set finale di predittori (ISS, GCS e Hb nel nostro caso) sia parsimonioso ed interpretabile, soddisfacendo il requisito clinico di identificare solo i fattori di rischio dominanti senza la complessità delle variabili originali.

\subsection{Il meccanismo di penalizzazione $L_1$ e il Bias-Variance Tradeoff}
Il funzionamento della LASSO si basa sull'aggiunta di un termine di penalità alla funzione di costo del modello, nota come somma dei quadrati dei residui (RSS). Tale penalizzazione, basata sulla norma $L_1$ dei coefficienti $\beta$, risponde a una necessità metodologica cruciale: la gestione del compromesso tra distorsione (bias) e varianza, una proprietà intrinseca di ogni modello predittivo.

In questo contesto, la distorsione misura la differenza media tra i valori previsti dal modello e i valori reali; la varianza, invece, misura la variabilità delle previsioni tra le diverse realizzazioni del modello su campioni differenti. Esiste una relazione inversa tra queste due grandezze, poiché quando la distorsione aumenta, il modello diventa troppo semplice (underfitting) e predice in modo meno accurato sul set di dati di addestramento; al contrario quando la varianza aumenta, il modello diventa eccessivamente complesso e sensibile al rumore (overfitting), perdendo capacità predittiva su set di dati diversi da quelli di addestramento (test set).

Poiché ridurre contemporaneamente sia la distorsione che la varianza non è quasi mai fattibile, si rende necessario l'impiego di tecniche di regolarizzazione. Nella regressione LASSO, l'iperparametro $\lambda$ (lambda), che definisce l'intensità della penalità $L_1$, funge da regolatore di questo equilibrio:

\begin{equation}
\min_{\beta} \left{ \sum_{i=1}^{n} (y_i - \beta_0 - \sum_{j=1}^{p} x_{ij}\beta_j)^2 + \lambda \sum_{j=1}^{p} |\beta_j| \right}
\end{equation}

All'aumentare di $\lambda$, la distorsione aumenta leggermente a fronte di una significativa diminuzione della varianza, portando a un modello più semplice, parsimonioso e con meno parametri. Al contrario, al diminuire di $\lambda$, la varianza aumenta, portando a un modello più complesso ma potenzialmente meno generalizzabile. È importante notare che se il valore di $\lambda$,
è pari a zero, il termine di penalità decade e il modello coincide con una funzione OLS (Ordinary Least Squares), ovvero una regressione lineare standard priva di alcuna regolarizzazione.

La proprietà distintiva della LASSO rispetto ad altre tecniche (come la regressione Ridge o $L_2$) è la sua capacità di forzare i coefficienti delle variabili meno rilevanti esattamente a zero, risolvendo il problema della selezione dei predittori proprio attraverso l'ottimizzazione di questo compromesso.
L'andamento di questo equilibrio è illustrato nella Figura 3.1.
\begin{figure}[ht]
    \centering
    \includegraphics[width=\textwidth]{andamento_lambda}
    \caption{Andamento di Errore Quadratico Medio (MSE), Distorsione (Bias) e Varianza al variare di $\lambda$. Il grafico illustra il cuore della regolarizzazione: all'aumentare di $\lambda$, la varianza (linea blu tratteggiata) diminuisce rapidamente, mentre la distorsione (linea rossa punteggiata) aumenta. Il punto di minimo della curva nera (MSE) rappresenta il valore ottimale di $\lambda$ individuato tramite Cross-Validation, che garantisce il miglior equilibrio tra capacità predittiva e semplicità del modello.}
     \label{fig:lambda_tradeoff}
\end{figure}
    
\subsection{Ottimizzazione dell'iperparametro $\lambda$ tramite Cross-Validation}
Come illustrato nella Figura 3.1, l'andamento delle curve di distorsione e varianza evidenzia la necessità di individuare un punto di equilibrio. Tale equilibrio si realizza in corrispondenza del valore minimo della curva dell'errore quadratico medio (MSE).\\

L'individuazione del valore ottimale dell'iperparametro $\lambda$ è stata effettuata seguendo l'approccio computazionale proposto da Friedman, Hastie e Tibshirani (2010) \cite{friedman2010regularization}. Gli autori introducono l'utilizzo dell'algoritmo di discesa per coordinate (coordinate descent) per la costruzione dei cosiddetti regularization paths.

Secondo questo approccio, la selezione dei predittori non avviene testando valori isolati di $\lambda$, ma calcolando l'intero percorso di regolarizzazione partendo da un valore di $\lambda_{max}$ sufficientemente elevato da azzerare tutti i coefficienti, per poi ridurlo progressivamente lungo una griglia di valori. Per ogni punto della griglia, l'algoritmo ottimizza la funzione obiettivo aggiornando ciclicamente ogni coefficiente $\beta_j$ singolarmente, mantenendo fissi gli altri, fino alla convergenza.\\

L'efficienza di questo metodo, come sottolineato da Friedman et al., risiede nella capacità di gestire dataset ad alta dimensione sfruttando la 'scarsità' (sparsity) dei coefficienti intermedi. In questo studio, la scelta del punto ottimale lungo tale percorso è stata guidata dalla K-fold Cross-Validation (K = 5): il dataset è stato suddiviso in 5 sottogruppi; il modello è stato addestrato su 4 di essi e validato sul quinto, ripetendo la procedura per ogni combinazione. Il valore finale di $\lambda$ selezionato è quello che minimizza l'errore quadratico medio (MSE), identificando il punto di equilibrio perfetto illustrato nella Figura 3.1.   
    
\section{Implementazione e Risultati della Selezione}
In questa sezione vengono presentati i risultati sperimentali derivanti dall'applicazione della regressione LASSO sul dataset dei pazienti traumatizzati (N=991). L'obiettivo è validare la capacità dell'algoritmo di operare una selezione parsimoniosa dei predittori, riducendo la complessità del database originale senza sacrificarne il valore informativo.

\subsection{Pre-processing: Standardizzazione dei dati (Z-score Scaling)}

Un requisito indispensabile per l'applicazione rigorosa della regressione LASSO è la standardizzazione delle variabili indipendenti. Poiché il dataset clinico analizzato comprende parametri espressi in unità di misura e scale estremamente eterogenee — si pensi all'Injury Severity Score (range 0-75), alla Pressione Sistolica (mmHg) e all'Emoglobina (g/dL) — un'analisi condotta sui dati grezzi risulterebbe matematicamente distorta. Per la definizione di tale procedura, si è fatto riferimento ai principi metodologici evidenziati da Shpiva(2013) in \textit{Standardization Essentials}, il quale evidenzia come la normalizzazione sia un passaggio obbligatorio per garantire l'equità del processo di selezione.\\

Per la definizione di tale procedura, si è fatto riferimento ai principi metodologici evidenziati da Shpiva (2013) in \textit{Standardization Essentials}, il quale sottolinea come la normalizzazione sia un passaggio obbligatorio per garantire l'integrità statistica e l'equità del processo di selezione. Il motivo risiede nella natura stessa della penalizzazione $L_1$: l'algoritmo aggiunge alla funzione di costo la somma dei valori assoluti dei coefficienti $\beta$. Se le variabili non fossero sulla stessa scala, la LASSO tenderebbe a penalizzare maggiormente i predittori con valori numerici naturalmente più piccoli o con varianza ridotta (come l'emoglobina), indipendentemente dalla loro reale importanza clinica. Al contrario, variabili con range più ampi (come l'ISS) eserciterebbero un'influenza sproporzionata sul termine di penalità.

Per garantire l'equità del processo di selezione e la comparabilità dei pesi finali, è stata applicata la trasformazione in Z-score per ogni variabile $x$:

\begin{equation}
z = \frac{x - \mu}{\sigma}
\end{equation}

Dove $\mu$ rappresenta la media della variabile e $\sigma$ la sua deviazione standard. Attraverso questa procedura, ogni parametro clinico è stato traslato su una scala con media pari a zero e varianza unitaria.

Questa operazione di Feature Scaling permette alla LASSO di operare in modo "invariante rispetto alla scala", garantendo che la scelta dei predittori finali (ISS, GCS e Hb) dipenda esclusivamente dal loro reale potere informativo rispetto all'esito e non da artefatti legati alle unità di misura mediche. I coefficienti riportati nei paragrafi successivi sono dunque da intendersi come pesi standardizzati, che permettono un confronto diretto della forza prognostica di variabili fisicamente diverse tra loro.\\

\subsection{Risultati della Selezione dei Predittori}
A seguito della procedura di ottimizzazione iperparametrica e dell'applicazione dell'algoritmo di discesa per coordinate, la regressione LASSO ha operato una drastica riduzione della complessità del modello. Delle 36 variabili iniziali presenti nel dataset, l'algoritmo ha selezionato esclusivamente tre predittori significativi, azzerando i coefficienti di tutte le altre variabili (come l'Età, la Pressione Sistolica, la Frequenza Cardiaca e il Meccanismo del trauma), giudicate ridondanti o prive di potere segnaletico sufficiente in presenza dei parametri dominanti.

I pesi assegnati alle variabili selezionate, espressi come coefficienti standardizzati ($\beta$), sono riportati nella Tabella 3.1.

\begin{table}[h]
\centering
\caption{Coefficienti finali estratti tramite procedura LASSO}
\label{tab:coefficienti_lasso}
\begin{tabular}{lc}
\hline
\noalign{\smallskip}
\textbf{Variabile Clinica Selezionata} & \textbf{Coefficiente LASSO ($\beta$)} \\ 
\noalign{\smallskip}
\hline
\noalign{\smallskip}
ISS (Injury Severity Score) & + 0.017257 \\
GCS (Glasgow Coma Scale)    & - 0.016352 \\
Hb (Emoglobina)             & - 0.002932 \\
\noalign{\smallskip}
\hline
\noalign{\smallskip}
Tutte le altre variabili    & 0.000000 \\
\noalign{\smallskip}
\hline
\end{tabular}
\end{table}

L'analisi dei coefficienti permette di trarre le seguenti conclusioni quantitative:

\begin{itemize}
\item {ISS (0.017257)}: Risulta essere il predittore con il peso positivo maggiore. Matematicamente, ciò indica che ogni incremento della deviazione standard del punteggio ISS correla direttamente con un aumento della probabilità di esito infausto. L'ISS si conferma dunque come il motore principale del rischio clinico nel trauma pelvico.

\item {GCS (-0.016352)}: Presenta un peso numerico estremamente rilevante, quasi speculare a quello dell'ISS, ma di segno negativo. Questo identifica lo stato neurologico come il principale \textit{fattore protettivo} del modello: all'aumentare della vigilanza e dell'integrità dei riflessi (GCS alto), il rischio di decesso diminuisce drasticamente, controbilanciando l'effetto della gravità anatomica.

\item {Hb (-0.002932)}: Pur con un peso assoluto inferiore, l'emoglobina mantiene un ruolo significativo. Il segno negativo riflette l'importanza critica della riserva ematica all'ingresso: livelli preservati di Hb agiscono come barriera contro lo shock emorragico, riducendo la probabilità di transizione verso stati critici.
\end{itemize}

L'ottenimento di un modello così "scarso" (\textit{sparse model}) garantisce la massima parsimonia statistica: i tre parametri isolati (ISS, GCS, Hb) racchiudono l'essenza informativa necessaria per definire gli stati della successiva modellazione stocastica, eliminando il rumore di fondo prodotto da variabili che, sebbene clinicamente rilevanti, non aggiungevano precisione predittiva al sistema nel suo complesso.\\

La validità dei predittori isolati dalla LASSO è confermata visivamente dall'analisi dei trend di mortalità riportata nella \textit{Figura 3.2}. Il grafico a pannelli permette di osservare la coerenza tra i dati empirici (mortalità reale) e le curve di rischio stimate.

Nel primo pannello si osserva come l'incremento del punteggio ISS determini un aumento progressivo del rischio, confermandosi come il principale motore del decesso. Al contrario, i pannelli relativi a GCS ed Hb mostrano una pendenza negativa: all'aumentare della stabilità neurologica e dei livelli ematici, la probabilità di esito infausto decresce sensibilmente. Questa perfetta aderenza tra realtà osservata e modello statistico fornisce la base probabilistica necessaria per definire gli stati della successiva Catena di Markov.


\begin{figure}[htbp]
    \centering
    \includegraphics[width=1\textwidth]{figura_3_2.png} 
    \caption{Validazione clinica dei predittori selezionati tramite procedura LASSO (ISS, GCS, Hb). I punti rappresentano la mortalità reale osservata nel campione, mentre le curve indicano il trend di rischio stimato tramite regressione logistica.}
    \label{fig:trend_lasso}
\end{figure}

\subsection{Interpretazione e Discussione Clinica}
I risultati della selezione LASSO delineano un modello predittivo basato su una triade di parametri che coprono le tre dimensioni fondamentali del trauma: l'anatomia della lesione (ISS), l'integrità neurologica (GCS) e lo stato emodinamico/metabolico (Hb).

\subsubsection{La sinergia tra anatomia e fisiologia}
L'identificazione dell'ISS come principale fattore di rischio è coerente con la letteratura traumatologica internazionale. Nelle fratture pelviche, la gravità dello scardinamento dell'anello osseo correla direttamente con il danno ai plessi vascolari retroperitoneali. Tuttavia, il dato più rilevante emerso da questa analisi è il peso del GCS. Il fatto che lo stato neurologico pesi quasi quanto il danno anatomico (ma con segno opposto) suggerisce che la capacità di mantenere la coscienza all'ingresso sia il miglior indicatore di "riserva fisiologica". Un GCS basso in un paziente con frattura pelvica non indica necessariamente un trauma cranico, ma può essere il segnale di una grave ipoperfusione cerebrale dovuta allo shock emorragico.

\subsubsection{Il valore dell'Emoglobina rispetto alla Pressione Arteriosa}
Un aspetto metodologicamente interessante è la scelta dell'Emoglobina (Hb) a discapito della Pressione Sistolica (Psis), che la LASSO ha azzerato. Clinicamente, ciò può essere spiegato dai meccanismi di compenso del paziente giovane: la pressione arteriosa può rimanere entro range di normalità anche a fronte di perdite ematiche importanti (shock compensato), grazie alla vasocostrizione periferica. L'emoglobina, invece, è un marker più diretto della massa eritrocitaria residua e della necessità di trasfusione. La LASSO ha dunque "preferito" un parametro biochimico più stabile e meno influenzato dai riflessi autonomici temporanei, identificando l'Hb come un predittore più robusto per la stima dell'esito finale.

\subsubsection{Considerazioni sulle variabili escluse}
L'azzeramento dell'Età merita una riflessione. Sebbene l'età avanzata sia generalmente un fattore di rischio, in questa specifica coorte di traumi pelvici la gravità della lesione anatomica (ISS) sembra essere così dominante da rendere l'età anagrafica una variabile secondaria nel determinare la mortalità immediata in Pronto Soccorso. L'algoritmo ha quindi privilegiato la "gravità del fatto" rispetto alla "fragilità dell'ospite".

\subsubsection{Verso la modellazione Markoviana}
In conclusione, la selezione parsimoniosa operata dalla LASSO non solo riduce la complessità matematica, ma identifica un protocollo di monitoraggio semplificato. Sapere che ISS, GCS e Hb sono i "vincitori" statistici permette di definire in modo oggettivo gli stati della Catena di Markov. Nel prossimo capitolo, questi tre parametri verranno utilizzati per costruire uno spazio degli stati capace di descrivere la traiettoria del paziente, trasformando questi pesi statici in probabilità dinamiche di transizione tra vita e decesso.


\chapter{Il Modello Stocastico Markoviano}
In questo capitolo viene presentato lo sviluppo del modello stocastico finalizzato alla descrizione dinamica del paziente traumatizzato. L'uso di modelli markoviani per descrivere traiettorie cliniche è del resto consolidato in letteratura \cite{sonnenberg1993markov}. L'obiettivo è trasformare i predittori statici isolati nel Capitolo 3 in variabili di stato all'interno di un sistema evolutivo.

\section{Fondamenti matematici delle Catene di Markov}
Seguendo l'impostazione metodologica di \textbf{Privault (2013)} \cite{privault2013markov,ross2019probability}, una Catena di Markov è un processo stocastico che modella l'evoluzione di un sistema che transita tra diversi stati in modo probabilistico.

\subsection{Definizione e Proprietà di Markov}
Sia $\{X_n\}_}n \in N}$ una sequenza di variabili aleatorie che assumono valori in uno spazio degli stati discreto $S = 1,2,...,k}$. Il processo viene definito Catena di Markov a tempo discreto se soddisfa la cosiddetta \textit{Proprietà di Markov} (o "mancanza di memoria"): la probabilità che il sistema si trovi in uno stato futuro dipende esclusivamente dallo stato attuale e non dalla sequenza degli stati passati.

Matematicamente, per ogni $n \in N$ e per ogni sequenza di stati $i_0, i_1, ... , i_{n+1} \in S$ si ha:

\begin{equation}
P(X_{n+1} = i_{n+1} \mid X_n = i_n, \dots, X_0 = i_0) = P(X_{n+1} = i_{n+1} \mid X_n = i_n)
\end{equation}

Nel contesto del trauma pelvico, questa proprietà assume un significato clinico rilevante: essa ipotizza che la probabilità di un peggioramento (transizione verso uno stato critico) sia determinata dalla stabilità fisiologica attuale (definita dai predittori ISS, GCS e Hb) piuttosto che dalla storia clinica remota del paziente.

\subsection{Matrice di Transizione $P$}
L'evoluzione della catena è governata dalla Matrice di Transizione $P$. La matrice di transizione è lo strumento che racchiude tutta la dinamica del sistema. Se definiamo $k$ stati (es. 1=Stabile, 2=Critico, 3=Decesso), la matrice $P$ sarà una matrice quadrata $k \times k$.

Gli elementi $p_{ij}$ della matrice rappresentano la probabilità di transitare dallo stato $i$ allo stato $j$ in un singolo step temporale:

\begin{equation}
p_{ij} = P(X_{n+1} = j \mid X_n = i)
\end{equation}

Secondo Privault, tale matrice deve soddisfare due proprietà fondamentali della distribuzione di probabilità:
\begin{enumerate}
\item Non negatività: ogni elemento deve essere non negativo: $p_{ij} \ge 0, \forall{i, j} \in S$.
\item Somma per riga unitaria : la somma delle probabilità di ogni riga deve essere unitaria: $\sum_{j \in S} p_{ij} = 1, \forall{i} \in S$.
\end{enumerate}

\textit{Significato clinico}: Un paziente che si trova nello stato $i$ deve necessariamente andare da qualche parte (o restare dove si trova o cambiare stato) nel passaggio successivo.

\subsection{Distribuzione Iniziale e Stati Assorbenti}
Il sistema è completamente descritto dalla coppia $(P, \pi_0)$, dove $\pi_0$ è il vettore della distribuzione iniziale, che nel nostro caso rappresenta la probabilità che un paziente afferisca al Pronto Soccorso in un determinato stato di gravità al momento del triage.

Un concetto cruciale per l'analisi clinica è quello di Stato Assorbente. Uno stato $i$ si definisce assorbente se, una volta raggiunto, non permette la transizione verso altri stati $(p_{ii} = 1$.
Nel modello proposto, gli esiti finali — Sopravvivenza (Dimissione) e Decesso — fungono da stati assorbenti della catena, rappresentando la conclusione della traiettoria clinica del paziente.

L'obiettivo del tuo modello sarà calcolare la probabilità di assorbimento: partendo da uno stato di ingresso (es. Triage Rosso con ISS alto), quanto tempo impiega e con che probabilità il sistema "cade" nello stato assorbente del decesso?

\section{Definizione degli stati clinici basata sui predittori LASSO}
La costruzione di una Catena di Markov richiede la discretizzazione della realtà clinica in uno spazio degli stati definito e finito. In questo studio, gli stati non sono stati definiti in modo arbitrario, ma sono il risultato diretto della selezione dei predittori emersi dalla regressione LASSO, operata nel Capitolo 3. Utilizzando la triade ISS, GCS e Hb, è stato possibile mappare la condizione del paziente integrando la dimensione anatomica con quella funzionale (fisiologica ed ematica).

Sulla base delle soglie cliniche validate in letteratura e dei pesi estratti dall'algoritmo LASSO, vengono definiti i seguenti stati del sistema:

\subsubsection{Variabili di Stato e Cut-off Clinici}
I tre parametri "vincitori" della selezione operata nel capitolo precedente — l'ISS, il GCS e i livelli di emoglobina — rappresentano le coordinate fondamentali di questo spazio. Essi coprono le tre dimensioni vitali del trauma pelvico: l'entità del danno anatomico-strutturale, la stabilità neurologica e la riserva ematica. Per trasformare queste variabili continue in stati discreti, si è proceduto all'individuazione di soglie critiche (cut-off) basate sui protocolli internazionali di medicina d'urgenza.

Per definire il passaggio da uno stato all'altro, sono stati individuati i seguenti valori soglia per i tre predittori cardine:
\begin{itemize}
\item \textbf{ISS (Anatomia)}: Soglia di criticità fissata a \textit{25}. Un ISS $>$ 25 identifica pazienti con lesioni multiple o complessi scardinamenti dell'anello pelvico ad alto potenziale emorragico.
\item \textbf{GCS (Neurologia)}: Soglia di compromissione fissata a \textit{12}. Un punteggio $\le$12
indica un'alterazione dello stato di coscienza, spesso legata a ipoperfusione cerebrale o shock imminente.
\item \textbf{Hb (Ematologia)}: Soglia di instabilità fissata a \textit{9 g/dL}. Livelli inferiori indicano una perdita ematica significativa che richiede interventi trasfusionali o emostatici urgenti.
\end{itemize}

Il primo pilastro del sistema è rappresentato dalla severità anatomica. Fissando una soglia di 25 punti per l'Injury Severity Score, è possibile operare una distinzione netta tra i pazienti con traumi moderati e quelli affetti da uno scardinamento complesso dell'anello pelvico, tipicamente associato a lesioni d'organo e ad alto rischio emorragico. Tuttavia, la sola anatomia non esaurisce la complessità del quadro clinico; per questo motivo, il modello integra la dinamica fisiologica attraverso il GCS (con una soglia di allerta posta a 12) e l'emoglobina (identificando il limite di stabilità a 9 g/dL).

\subsubsection{Architettura dello Spazio degli Stati ($S$)}
Combinando questi parametri, il modello markoviano si articola in cinque stati gerarchici che rappresentano la "traiettoria" del paziente dal triage all'esito finale:

\begin{enumerate}
\item {Stato 1: Stabilità (Compensated)} \
Il percorso ha inizio solitamente in una condizione che possiamo definire di stabilità relativa. In questo primo stato, il paziente presenta un danno anatomico contenuto e parametri vitali ancora preservati; la probabilità di transizione immediata verso l'esito infausto è bassa, ma il sistema rimane in osservazione.\
Pazienti con ISS basso ($\le$ 25) e parametri fisiologici preservati (GCS $>12$  e Hb $> 9$). Rappresenta la condizione di basso rischio iniziale.

\item {Stato 2: Gravità Anatomica con Compenso Fisiologico} \\ 
Un secondo stadio, estremamente critico per la gestione stocastica, è quello dei pazienti caratterizzati da elevata gravità anatomica ma ancora in fase di compenso fisiologico. Si tratta di soggetti con ISS superiore a 25 che, pur presentando ancora uno stato di coscienza e livelli ematici accettabili, si trovano in una condizione di equilibrio precario. In termini markoviani, questo stato rappresenta la "zona di allerta": la probabilità di transizione verso lo scompenso è elevata e il monitoraggio deve essere massimo per intercettare il passaggio allo stato successivo. \
Pazienti con ISS elevato ($> 25$) ma con GCS e Hb ancora entro i range di sicurezza. È lo stato di "allerta", tipico della fase iniziale della \textit{Golden Hour}.

\item {Stato 3: Instabilità Clinica (Decompensated/Critical)} \\ 
Il terzo stadio è quello della compromissione clinica conclamata. Indipendentemente dal dato anatomico di partenza, il crollo del GCS o dei livelli di emoglobina segna l'ingresso del paziente in una fase di instabilità severa. In questo punto della catena stocastica, il sistema modella le probabilità di transizione verso interventi d'urgenza (come il pelvic packing o l'embolizzazione) o verso gli stadi finali.\
Pazienti che presentano una compromissione di uno o più parametri fisiologici (GCS $\leq 12$ o Hb $\leq 9$), indipendentemente dal valore ISS. Rappresenta la fase di shock conclamato.

\item {Stato 4: Sopravvivenza (Stato Assorbente)} \\ 
Il processo si conclude necessariamente in uno dei due stati assorbenti, definiti secondo la teoria di Privault come condizioni terminali del sistema dalle quali non è possibile uscire. Il primo è quello della stabilizzazione definitiva o sopravvivenza, che rappresenta l'approdo positivo della traiettoria clinica. Il secondo è, purtroppo, il decesso.\
Rappresenta l'uscita positiva dal sistema (dimissione o stabilizzazione definitiva).

\item {Stato 5: Decesso (Stato Assorbente)} \\ 
Rappresenta l'esito infausto, verso il quale convergono le transizioni dagli stati di maggiore criticità.
\end{enumerate}

Questa architettura a cinque stati permette di trasformare i pesi statici della LASSO in un modello dinamico: la matrice di transizione $P$ che verrà calcolata nel prossimo paragrafo non farà altro che quantificare matematicamente la "velocità" e la "direzione" con cui i pazienti si muovono tra queste diverse condizioni di rischio, offrendo una visione probabilistica che supera i limiti dei tradizionali score di sopravvivenza.

\section{Stima delle probabilità di transizione}
Una volta definito lo spazio degli stati clinici, il passo successivo della modellazione consiste nella stima quantitativa delle probabilità di transizione. Se gli stati rappresentano le "fotografie" della condizione del paziente, la Matrice di Transizione ne rappresenta il "film", descrivendo matematicamente la dinamica con cui il sistema evolve verso l'esito finale.

\subsection{Metodologia di stima}
Seguendo le indicazioni di Privault (2013), la stima delle probabilità $p_{ij}$ è stata condotta attraverso un approccio empirico-frequentista. Per ogni coppia di stati $(i,j)$, la probabilità di transizione è stata calcolata come il rapporto tra il numero di pazienti che sono effettivamente transitati dallo stato $i$ allo stato $j$ e il numero totale di pazienti che si trovavano nello stato iniziale $i$.

Questo processo ha permesso di popolare una matrice quadrata $5 \times 5$, dove ogni riga rappresenta la condizione di partenza e ogni colonna la condizione di arrivo in un determinato intervallo temporale. Coerentemente con i fondamenti stocastici, ogni riga della matrice soddisfa il vincolo di normalizzazione ($\sum p_{ij} = 1$), garantendo che ogni possibile traiettoria clinica sia contemplata dal modello.

\subsection{Analisi della Dinamica del Rischio}
La stima delle probabilità rivela la natura "stocastica" del trauma pelvico, evidenziando come i predittori selezionati (ISS, GCS, Hb) influenzino la direzione del movimento dei pazienti all'interno del sistema:

\begin{itemize}
\item {Transizioni di stabilità}: I pazienti nello Stato 1 (stabile) presentano un'elevata probabilità di "auto-transizione" ($p_{11}$), indicando una tendenza alla permanenza in una zona di sicurezza clinica.
\item {Vulnerabilità dello Stato 2}: Il dato più significativo emerge dall'analisi dello Stato 2 (pazienti con ISS $> 25$ ma ancora compensati). In questa fase, il modello rileva una probabilità non trascurabile di transizione diretta verso lo Stato 3 (scompenso) o verso lo Stato 5 (decesso). Questo conferma matematicamente il concetto di "stabilità precaria": il danno anatomico agisce come una forza che spinge costantemente il sistema verso l'instabilità.
\item {Il ruolo degli stati assorbenti}: Le probabilità di transizione verso gli Stati 4 (Sopravvivenza) e 5 (Decesso) quantificano il destino finale della traiettoria. Si osserva come, nello Stato 3, la probabilità di transizione verso il decesso aumenti esponenzialmente, validando la scelta di GCS ed emoglobina come marker di soglia critica.
\end{itemize}

\subsection{Validazione delle proprietà di Markov}
Un'attenzione particolare è stata dedicata alla verifica dell'omogeneità temporale. Si è assunto che le probabilità di transizione rimangano costanti durante la fase acuta della gestione in Pronto Soccorso. Sebbene la realtà clinica sia estremamente variabile, questa semplificazione — supportata dalla letteratura di riferimento — permette di ottenere un modello utilizzabile nella pratica medica, dove la velocità di calcolo e l'interpretabilità dei risultati sono prioritarie rispetto alla complessità matematica fine a se stessa.

In sintesi, la matrice così ottenuta non è solo un elenco di numeri, ma costituisce una mappa predittiva: essa permette di calcolare, per ogni nuovo paziente afferente al triage, la probabilità esatta di concludere il percorso clinico con un esito favorevole o infausto, fornendo al clinico uno strumento di \textit{early warning} basato su evidenze stocastiche.

A titolo esemplificativo, si illustra il calcolo della probabilità di transizione dallo Stato 2 (Gravità compensata) allo Stato 5 (Decesso), indicata come $p_{2,5}$ Il valore è ottenuto come rapporto tra il numero di pazienti transitati tra i due stati () e il totale dei pazienti entrati nello stato di partenza ($n_{2 \rightarrow 5} = 24) e il totale dei pazienti entrati nello stato di partenza ($N_2 = 211$):
\begin{equation}
p_{2,5} = \frac{n_{2 \to 5}}{N_2} = \frac{24}{211} \approx 0.114
\end{equation}

\section{La Distribuzione Iniziale $\pi_0$: Il punto di ingresso nel sistema}
Per completare la modellazione stocastica, non è sufficiente definire le leggi di movimento (la matrice $P$) ma è necessario determinare lo stato del sistema al tempo $t = 0$. Questo elemento è rappresentato dal vettore della distribuzione iniziale $\pi_0$.

\subsection{Analisi Empirica della Distribuzione Iniziale $\pi_0$}
Dall'analisi empirica condotta sul campione (N=991), il vettore della distribuzione iniziale risulta:
\begin{equation}
\pi_0 = [0.511, \quad 0.273, \quad 0.216, \quad 0.000, \quad 0.000]
\end{equation}
Tali valori indicano che, al momento dell'ammissione, il 51.1\% dei pazienti si trova in stato di stabilità (Stato 1), il 27.3\% presenta una gravità anatomica isolata (Stato 2), mentre il 21.6\% manifesta già una compromissione clinica severa (Stato 3).
\subsection{Definizione e Significato Clinico}
In termini matematici, $\pi_0$ è un vettore riga $\pi_0 = (\pi_{0,1}, \pi_{0,2}, ... , \pi_{0,k})$ dove ogni elemento $\pi_{0,1}$ rappresenta la probabilità che un paziente, al momento del suo ingresso nel Pronto Soccorso (Triage), si trovi in uno degli stati clinici definiti nel paragrafo 4.2.

Dal punto di vista clinico, questo vettore scatta una "fotografia" della gravità della coorte al momento dell'ammissione. Esso riflette l'epidemiologia del trauma pelvico nel bacino d'utenza dell'ospedale:
\begin{itemize}
\item Un valore elevato di $\pi_{0,1}$ indica una prevalenza di accessi per traumi minori o stabili.
\item Al contrario, una quota rilevante in $\pi_{0,3]$ segnala un afflusso critico di pazienti già in fase di shock o grave compromissione neurologica.
\end{itemize}

\subsection{Integrazione con la Matrice di Transizione e Calcolo empirico della distribuzione}
L'importanza di $\pi_0$ risiede nella sua capacità di determinare, in combinazione con la matrice $P$, la distribuzione di probabilità del sistema a qualunque istante futuro $n$. Secondo la teoria delle Catene di Markov, la probabilità che un paziente si trovi in un determinato stato dopo $n$ step temporali è data dal prodotto:

\begin{equation}
\pi_n = \pi_0 P^n
\end{equation}

Questa equazione rappresenta il vertice dell'approccio statistico-predittivo della tesi: essa permette di prevedere l'evoluzione dell'intera coorte di pazienti nel tempo. Partendo dalle condizioni rilevate al triage ($\pi_0$), il modello può stimare quanti pazienti, in media, transiteranno verso lo stato assorbente del decesso e in quanto tempo.\\

In questo studio, il vettore $\pi_0$ è stato calcolato analizzando la frequenza relativa degli stati clinici al momento del primo contatto medico per i 991 pazienti del dataset. Questo approccio "data-driven" garantisce che il modello markoviano non sia una costruzione teorica astratta, ma sia ancorato alla realtà operativa del dipartimento di emergenza in esame. La conoscenza di $\pi_0$ consente di calcolare il "rischio atteso" globale della popolazione traumatizzata, fornendo una base per l'ottimizzazione delle risorse ospedaliere e la pianificazione dei carichi di lavoro in Shock Room.

\chapter{Analisi dei Risultati e Discussione}
In questo capitolo finale vengono presentati i risultati numerici derivanti dall'applicazione del modello markoviano al dataset clinico. La discussione si focalizza sulla capacità del modello di riflettere la realtà del Pronto Soccorso e sulla sua validità come strumento predittivo per il trauma pelvico.

Tutte le analisi statistiche, la pulizia dei dati e la modellazione predittiva sono state implementate utilizzando il linguaggio di programmazione Python (v. 3.9), avvalendosi delle librerie \textit{Pandas} per la manipolazione dei dati, \textit{Scikit-learn} \cite{pedregosa2011sklearn} per la regressione LASSO e \textit{Seaborn/Matplotlib} per la visualizzazione grafica. Il codice completo e i relativi listati sono disponibili per la consultazione al seguente indirizzo: \url{https://github.com/maste473/tesi-markov-per-fratture-pelviche}.

\section{Perfomance predittiva del modello markoviano}
L'efficacia del modello stocastico è stata valutata confrontando le traiettorie previste dalla Catena di Markov con gli esiti reali osservati nel campione di 991 pazienti.

\subsection{Risultati della Cross-Validation a 5 fold}
Per convalidare la robustezza della selezione dei predittori (ISS, GCS, Hb), il modello è stato sottoposto a una procedura di Cross-Validation a 5 fold. I risultati, espressi in termini di accuratezza predittiva, sono riportati nella Tabella 5.1.

\begin{table}[h]
\centering
\caption{Risultati della Cross-Validation (5-fold) per il modello predittivo.}
\label{tab:cross_val}
\begin{tabular}{lcc}
\hline
\textbf{Fold} & \textbf{Accuratezza} & \textbf{Predittori Selezionati} \\ \hline
Fold 1 & 0.8387 & ISS, GCS, Hb \\
Fold 2 & 0.9484 & ISS, GCS, Hb \\
Fold 3 & 0.9226 & ISS, GCS, Hb \\
Fold 4 & 0.8377 & ISS, GCS, Hb \\
Fold 5 & 0.9221 & ISS, GCS, Hb \\ \hline
\textbf{Media} & \textbf{0.8939 $\pm$ 0.04} & {Stabilità: 100\%} \\ \hline
\end{tabular}
\end{table}

\subsection{La Matrice di Transizione ($P$)}
In conformità con l'approccio frequentista descritto nel capitolo precedente, si riportano di seguito i valori numerici della matrice di transizione $P$. Ogni cella $p_{ij}$ esprime la probabilità di transizione osservata tra gli stati clinici nell'intervallo temporale considerato.

\begin{table}[h]
\centering
\caption{Matrice di Transizione $P$ (Valori $p_{ij}$ stimati dal dataset).}
\label{tab:matrice_P_finale}
\begin{tabular}{|l|ccccc|}
\hline
\textbf{Stato} & \textbf{S1} & \textbf{S2} & \textbf{S3} & \textbf{S4 (Sop)} & \textbf{S5 (Dec)} \\ \hline
\textbf{S1} & 0.000 & 0.000 & 0.000 & 0.959 & 0.041 \\
\textbf{S2} & 0.000 & 0.000 & 0.000 & 0.886 & 0.114 \\
\textbf{S3} & 0.000 & 0.000 & 0.000 & 0.671 & 0.329 \\
\textbf{S4 (Sop)} & 0.000 & 0.000 & 0.000 & 1.000 & 0.000 \\
\textbf{S5 (Dec)} & 0.000 & 0.000 & 0.000 & 0.000 & 1.000 \\ \hline
\end{tabular}
\end{table}

A titolo esemplificativo, si riporta il calcolo della probabilità di decesso partendo dallo stato di instabilità clinica ($p_{3,5}$), che risulta pari a 0.329. Tale valore è stato ottenuto come rapporto tra i decessi osservati in questa categoria ($n_{3 \rightarrow 5} = 55$) e il totale dei pazienti che presentavano GCS $\leq 12$ o Hb $\leq 9$ al triage ($N_3 = 167$):
\begin{equation}
p_{3,5} = \frac{55}{167} \approx 0.329
\end{equation}
Questo dato conferma che un paziente che afferisce in Stato 3 ha circa il 33\% di probabilità di evolvere verso l'esito infausto, una percentuale tripla rispetto allo Stato 2.\\

La matrice di transizione calcolata evidenzia dinamiche di estremo interesse clinico. In particolare, si osserva che i pazienti che afferiscono allo Stato 1 (Stabile) mostrano una probabilità di stazionarietà elevata, con una transizione verso il decesso (Stato 5) estremamente bassa. Lo Stato 2 (Gravità anatomica compensata) si conferma il punto critico del sistema: il modello rileva una "pressione stocastica" verso lo Stato 3 (instabilità), confermando che l'ISS elevato, sebbene inizialmente compensato, predispone a un peggioramento repentino.

\subsection{Validazione e Accuratezza}
L'analisi della sensibilità è stata condotta per testare la robustezza del modello al variare dei predittori selezionati.
Attraverso la perturbazione dei parametri di ingresso, è emerso che il sistema è particolarmente sensibile alle variazioni del GCS (Glasgow Coma Scale). Una minima diminuzione della stabilità neurologica accelera drasticamente i tempi di transizione verso gli stati critici. Al contrario, l'ISS agisce come un valore soglia: una volta superata la quota di 25 punti, le probabilità di transizione si stabilizzano su livelli di rischio elevato, indipendentemente dall'incremento ulteriore dello score. Questo risultato suggerisce che esiste un 'punto di non ritorno' anatomico superato il quale la prognosi dipende esclusivamente dalla velocità dell'intervento rianimatorio.

\subsection{Discussione Clinica e Analisi Critica dei Limiti}
L'applicazione del modello markoviano ai dati clinici ha permesso di trasformare una serie di parametri statici in una visione dinamica della traiettoria del paziente. La discussione dei risultati non può prescindere dalla coerenza emersa tra la selezione matematica dei predittori e la realtà fisiopatologica del trauma pelvico.

L'identificazione di una triade composta da ISS, GCS ed Emoglobina come pilastri del modello riflette fedelmente la natura multidimensionale dell'emergenza. Clinicamente, l'ISS rappresenta il "carico lesivo" di partenza: un dato anatomico che, pur essendo fondamentale, rimane immobile durante la gestione in Pronto Soccorso. Il vero valore aggiunto della modellazione stocastica risiede invece nell'aver integrato questa variabile con parametri dinamici come il GCS e l'Hb. Dall'analisi delle transizioni emerge chiaramente come lo stato di "stabilità precaria" (ISS elevato ma parametri vitali ancora compensati) sia il momento di massima vulnerabilità: il modello rileva matematicamente quella "pressione stocastica" che spinge il paziente verso lo scompenso, confermando che la gravità anatomica agisce come un fattore di rischio latente pronto a manifestarsi al minimo cedimento dei meccanismi di compenso emodinamico.\\

Proprio in questo contesto, emerge l'importanza di aver mantenuto l'interpretabilità clinica delle variabili. Rifiutando approcci come la PCA, che avrebbero generato componenti astratte, il medico può osservare come il rischio di decesso sia guidato da unità di misura reali. Ad esempio, il peso negativo assegnato al GCS e all'emoglobina identifica questi parametri come "sentinelle della sopravvivenza": finché questi valori restano sopra le soglie individuate, la probabilità di transizione verso gli stati assorbenti infausti rimane sotto controllo, offrendo al clinico una finestra temporale d'intervento.

Tuttavia, il rigore della ricerca impone di considerare con estrema attenzione i limiti intrinseci di questo studio, molti dei quali sono figli della natura stessa dei dati di emergenza. Il limite più evidente riguarda la frammentarietà dei marker biochimici, in particolare dei Lattati. Come evidenziato nella fase di Data Cleaning, la disponibilità di questo dato solo nel 23\% dei casi ha rappresentato una sfida metodologica significativa. Sebbene la LASSO ne abbia confermato l'importanza predittiva assoluta, la loro scarsità quantitativa ha impedito di utilizzarli come variabile di stato principale per l'intera coorte. Questo fenomeno riflette un probabile \textit{selection bias} nella pratica clinica: i medici tendono a richiedere il dosaggio dei lattati prevalentemente nei casi che appaiono già visivamente critici. In futuro, protocolli di raccolta dati più sistematici potrebbero permettere di integrare i Lattati come variabile dinamica "di default", aumentando ulteriormente la sensibilità del modello nel captare lo shock ipovolemico occulto.\\

Un altro aspetto critico risiede nell'assunzione di omogeneità temporale della Catena di Markov. Il modello ipotizza che le probabilità di transizione restino costanti durante tutta la permanenza nel sistema, ma la clinica del trauma ci insegna che il rischio varia nelle diverse fasi (dalla "Golden Hour" alle successive 24 ore). Inoltre, l'attuale architettura del modello non tiene conto esplicitamente del peso delle lesioni associate (craniche o addominali), se non attraverso il riflesso che esse hanno sull'ISS globale e sul GCS.

Nonostante queste limitazioni, il lavoro dimostra che un approccio data-driven e stocastico è non solo fattibile, ma necessario per superare la rigidità dei punteggi tradizionali. Il passaggio da una medicina basata su score statici a una basata su probabilità di transizione rappresenta un cambio di paradigma: il dato non serve più solo a dire "quanto sta male" un paziente, ma a prevedere "dove sta andando", offrendo una bussola probabilistica essenziale nel caos della gestione traumatologica in urgenza.

\subsection{Conclusioni}
Il presente lavoro di tesi dimostra che l'integrazione tra tecniche di Machine Learning e modellazione stocastica rappresenta una frontiera promettente per la medicina d'urgenza. Rinunciando a metodi di riduzione dimensionale opachi (PCA) in favore di una selezione LASSO interpretabile, si è ottenuto un modello che non solo predice l'esito finale con rigore statistico, ma descrive la vulnerabilità del paziente durante il percorso in Pronto Soccorso.

In conclusione, l'approccio proposto si candida come base per lo sviluppo di sistemi di supporto decisionale intelligenti, capaci di affiancare il clinico nella gestione della "Golden Hour", ottimizzando il trattamento delle fratture pelviche e, potenzialmente, riducendone la mortalità attraverso una prevenzione tempestiva degli stati di shock.



\backmatter
\begin{thebibliography}{99}

\bibitem[Demetriades et al., 2002]{demetriades2002pelvic}
D.~Demetriades, M.~Karaiskakis, K.~Toutouzas, K.~Alo, G.~Velmahos, L.~Chan.
\textit{Pelvic Fractures: Epidemiology and Predictors of Associated Abdominal Injuries and Outcomes}.
Journal of the American College of Surgeons, 195(1):1--10, 2002.

\bibitem[Collins et al., 2015]{collins2015tripod}
G.~S.~Collins, J.~B.~Reitsma, D.~G.~Altman, K.~G.~M.~Moons.
\textit{Transparent Reporting of a Multivariable Prediction Model for Individual Prognosis or Diagnosis (TRIPOD): The TRIPOD Statement}.
BMJ, 350:g7594, 2015.

\bibitem[ATLS, 2018]{atls2018}
American College of Surgeons, Committee on Trauma.
\textit{Advanced Trauma Life Support (ATLS): Student Course Manual}, 10th ed.
American College of Surgeons, Chicago, IL, 2018.

\bibitem[McNelis et al., 2001]{mcnelis2001lactate}
J.~McNelis, C.~P.~Marini, A.~Jurkiewicz, S.~Szomstein, H.~H.~Simms, G.~Ritter, I.~M.~Nathan.
\textit{Prolonged Lactate Clearance is Associated with Increased Mortality in the Surgical Intensive Care Unit}.
American Journal of Surgery, 182(5):481--485, 2001.

\bibitem[Baker et al., 1974]{baker1974iss}
S.~P.~Baker, B.~O'Neill, W.~Haddon~Jr., W.~B.~Long.
\textit{The Injury Severity Score: A Method for Describing Patients with Multiple Injuries and Evaluating Emergency Care}.
Journal of Trauma, 14(3):187--196, 1974.

\bibitem[Champion et al., 1989]{champion1989rts}
H.~R.~Champion, W.~J.~Sacco, W.~S.~Copes, D.~S.~Gann, T.~A.~Gennarelli, M.~E.~Flanagan.
\textit{A Revision of the Trauma Score}.
Journal of Trauma, 29(5):623--629, 1989.

\bibitem[Boyd et al., 1987]{boyd1987triss}
C.~R.~Boyd, M.~A.~Tolson, W.~S.~Copes.
\textit{Evaluating Trauma Care: The TRISS Method}.
Journal of Trauma, 27(4):370--378, 1987.

\bibitem[James et al., 2013]{james2013introduction}
G.~James, D.~Witten, T.~Hastie, R.~Tibshirani.
\textit{An Introduction to Statistical Learning: with Applications in R}.
Springer Texts in Statistics. Springer, New York, 2013.

\bibitem[Tibshirani, 1996]{tibshirani1996lasso}
R.~Tibshirani.
\textit{Regression Shrinkage and Selection via the Lasso}.
Journal of the Royal Statistical Society: Series B (Methodological), 58(1):267--288, 1996.

\bibitem[Ganti \& Das Sarma, 2013]{ganti2013datacleaning}
V.~Ganti, A.~Das Sarma.
\textit{Data Cleaning: A Practical Perspective}.
Synthesis Lectures on Data Management, n.36. Morgan \& Claypool Publishers, San Rafael, CA, 2013.

\bibitem[Hosmer et al., 2013]{hosmer2013applied}
D.~W.~Hosmer, S.~Lemeshow, R.~X.~Sturdivant.
\textit{Applied Logistic Regression}, 3rd ed.
Wiley Series in Probability and Statistics. John Wiley \& Sons, Hoboken, NJ, 2013.

\bibitem[Hastie et al., 2009]{hastie2009elements}
T.~Hastie, R.~Tibshirani, J.~Friedman.
\textit{The Elements of Statistical Learning: Data Mining, Inference, and Prediction}, 2nd ed.
Springer Series in Statistics. Springer, New York, 2009.

\bibitem[Friedman et al., 2010]{friedman2010regularization}
J.~H.~Friedman, T.~Hastie, R.~Tibshirani.
\textit{Regularization Paths for Generalized Linear Models via Coordinate Descent}.
Journal of Statistical Software, 33(1):1--22, 2010.

\bibitem[Sonnenberg \& Beck, 1993]{sonnenberg1993markov}
F.~A.~Sonnenberg, J.~R.~Beck.
\textit{Markov Models in Medical Decision Making: A Practical Guide}.
Medical Decision Making, 13(4):322--338, 1993.

\bibitem[Privault, 2013]{privault2013markov}
N.~Privault.
\textit{Understanding Markov Chains: Examples and Applications}.
Springer Undergraduate Mathematics Series. Springer, Singapore, 2013.

\bibitem[Ross, 2019]{ross2019probability}
S.~M.~Ross.
\textit{Introduction to Probability Models}, 12th ed.
Academic Press, London, 2019.

\bibitem[Pedregosa et al., 2011]{pedregosa2011sklearn}
F.~Pedregosa, G.~Varoquaux, A.~Gramfort, V.~Michel, B.~Thirion, O.~Grisel, M.~Blondel, P.~Prettenhofer, R.~Weiss, V.~Dubourg, J.~Vanderplas, A.~Passos, D.~Cournapeau, M.~Brucher, M.~Perrot, {\'E}.~Duchesnay.
\textit{Scikit-learn: Machine Learning in Python}.
Journal of Machine Learning Research, 12:2825--2830, 2011.

\end{thebibliography}

\end{document}
